% ---
% id: receiver/detector-imperfections/lem:pure-loss-lift
% kind: lem
% title: "The pure-loss channel on the lift of a contraction"
% summary: >-
%   The adjoint of the pure-loss channel of transmittance t takes the lift of
%   a contraction X to the lift of tX + (1 − t)·1. Every joint no-click
%   element is such a lift, so loss in front of a receiver acts on its Grams
%   as a scalar.
% prominence: technical
% status: accepted
% uses: [def:pure-loss, receiver/lem:fock-maps]
% mentions: [prop:efficiency-normalization, receiver/thm:click-povm]
% assumes: []
% refs: []
% provenance: newly-derived
% lean:
% ---

The \tref{joint no-click elements} of a receiver are lifts $\Fock(\id-A_J)$ of
contractions (Theorem~\ref{thm:click-povm}\,(ii)), so the following identity
is what Proposition~\ref{prop:efficiency-normalization} needs from the
pure-loss channel.

\begin{lemma}\label{lem:pure-loss-lift}
Let $t\in[0,1]$ and let $X$ be an operator on $\mathbb{C}^{\InpC}$ with
$X^{\dagger}X\preceq\id$. Then
\begin{equation}\label{eq:purelosslift}
  (\Lambda^{\mathrm{loss}}_t)^{\dagger}\big(\Fock(X)\big)
  \;=\;\Fock\big(t\,X+(1-t)\,\id\big).
\end{equation}
\end{lemma}

\begin{proof}
\emph{Step 1 (the adjoint).} By Lemma~\ref{lem:fock-maps}\,(iii) with
$X_2=\id$, under the identification of Definition~\ref{def:pure-loss},
$\Fock(X)\otimes\idon{\mathcal{F}(\mathbb{C}^{\InpC'})}=\Fock(X\oplus\id)$,
an operator of norm at most one (Lemma~\ref{lem:fock-maps}\,(i)). So
\eqref{eq:purelossadj} and Lemma~\ref{lem:fock-maps}\,(i) give
\begin{equation*}
  (\Lambda^{\mathrm{loss}}_t)^{\dagger}\big(\Fock(X)\big)
  =\Fock(w_t)^{\dagger}\,\Fock(X\oplus\id)\,\Fock(w_t)
  =\Fock\big(w_t^{\dagger}(X\oplus\id)w_t\big).
\end{equation*}

\emph{Step 2 (the compression).} For $v\in\mathbb{C}^{\InpC}$ write $v'$ for
its copy in $\mathbb{C}^{\InpC'}$. Then $w_tv=\sqrt t\,v+\sqrt{1-t}\,v'$,
$(X\oplus\id)w_tv=\sqrt t\,Xv+\sqrt{1-t}\,v'$, and $w_t^{\dagger}$ maps
$a+b'$ to $\sqrt t\,a+\sqrt{1-t}\,b$, so
$w_t^{\dagger}(X\oplus\id)w_t=t\,X+(1-t)\,\id$.
\end{proof}
